\chapter{Preliminary Study: From the Starting Method to the Proposed Pipeline}
\thispagestyle{empty}

Before the proposed pipeline was designed, a broad set of representations and classifiers was
evaluated under the protocol of Chapter~4, starting from the method with the highest reported
accuracy and changing one part at a time. This chapter summarises that study in the order in
which the steps were taken; it explains why the proposed pipeline combines the decomposition of
the starting method with Riemannian geometry.

% result of one step on unseen people, from the generated numbers (letters A--J = steps 1--10)
\newcommand{\stepres}[1]{\par\smallskip\noindent\hspace*{1.5em}$\rightarrow$\ \textit{Unseen
  people:} DS-2 \textbf{\R{Step#1DStwoBal}\%} balanced accuracy, \R{Step#1DStwoSubj} children
  correct; DS-1 rest \R{Step#1DSoneRestBal}\%, \R{Step#1DSoneRestSubj} participants correct.\par
  \medskip}

\section{The steps}
\begin{description}[leftmargin=0pt, itemsep=2pt]
\item[Step 1: Starting method.] SMVMD of every window, nine statistical features per mode and
  channel combined by an energy-weighted average, mRMR feature selection~\cite{peng2005} and a
  $k$-nearest-neighbour classifier~\cite{cover1967}, with the settings of its
  description~\cite{chandela2025} (remaining details in Appendix~A).
  \stepres{A}
  This is far below the published figures (Section~\ref{sec:whysubject}). The features identify
  the person: a nearest-neighbour classifier recognises \emph{which} DS-1 participant a window
  comes from with \R{DSoneRestProbeAcc}\% accuracy (chance \R{DSoneRestProbeChance}\%).

\item[Step 2: Tuned settings.] The same features, with the number of mRMR features,
  standardisation, $k$, distance and neighbour weighting chosen by nested cross-validation.
  \stepres{B}

\item[Step 3: Spectral features.] Absolute and relative band powers, theta/beta ratio,
  spectral entropy and further time-domain measures, with logistic regression, SVM, random
  forest and XGBoost.
  \stepres{C}

\item[Step 4: Connectivity.] Coherence and weighted phase-lag index between all pairs of
  channels, with the same classifiers.
  \stepres{D}

\item[Step 5: Other signal-processing features.] Filter-bank log-variance, wavelet-packet
  energies and aperiodic ($1/f$) features~\cite{donoghue2020}.
  \stepres{E}

\item[Step 6: One vector per person.] The mean and standard deviation of each person's spectral
  features over their windows, classified at person level.
  \stepres{F}

\item[Step 7: Riemannian geometry.] The channel covariance matrix of each window, mapped to the
  Riemannian tangent space and classified with logistic regression or SVM.
  \stepres{G}

\item[Step 8: Convolutional networks.] EEGNet and ShallowConvNet trained on the raw windows.
  \stepres{H}

\item[Step 9: MiniRocket.] Random convolutional features~\cite{dempster2021} with a ridge
  classifier.
  \stepres{I}

\item[Step 10: Ensemble.] The mean probability of five earlier models, fixed in advance.
  \stepres{J}
\end{description}

Table~\ref{tab:progression} collects the results of all steps, and Fig.~\ref{fig:progression}
shows the progression on DS-2 up to the proposed pipeline of Chapter~6.

\begin{table}[H]
\caption{Preliminary study: best result of each step on unseen people (balanced window
accuracy, \%, and people classified correctly by majority vote), followed by the proposed
pipeline.}
\label{tab:progression}\centering\small
\begin{tabular}{clcccc}
\hline
Step & Method & \multicolumn{2}{c}{DS-1 rest} & \multicolumn{2}{c}{DS-2 ADHD$^\ast$} \\
 & & Bal.\ acc. & Participants & Bal.\ acc. & Children \\
\hline
1 & Starting method (SMVMD, mRMR, KNN) & \R{StepADSoneRestBal} & \R{StepADSoneRestSubj} & \R{StepADStwoBal} & \R{StepADStwoSubj} \\
2 & Tuned settings & \R{StepBDSoneRestBal} & \R{StepBDSoneRestSubj} & \R{StepBDStwoBal} & \R{StepBDStwoSubj} \\
3 & Spectral features & \R{StepCDSoneRestBal} & \R{StepCDSoneRestSubj} & \R{StepCDStwoBal} & \R{StepCDStwoSubj} \\
4 & Connectivity & \R{StepDDSoneRestBal} & \R{StepDDSoneRestSubj} & \R{StepDDStwoBal} & \R{StepDDStwoSubj} \\
5 & Filter bank, wavelet packet, aperiodic & \R{StepEDSoneRestBal} & \R{StepEDSoneRestSubj} & \R{StepEDStwoBal} & \R{StepEDStwoSubj} \\
6 & One vector per person & \R{StepFDSoneRestBal} & \R{StepFDSoneRestSubj} & \R{StepFDStwoBal} & \R{StepFDStwoSubj} \\
7 & Riemannian tangent space & \R{StepGDSoneRestBal} & \R{StepGDSoneRestSubj} & \R{StepGDStwoBal} & \R{StepGDStwoSubj} \\
8 & EEGNet, ShallowConvNet & \R{StepHDSoneRestBal} & \R{StepHDSoneRestSubj} & \R{StepHDStwoBal} & \R{StepHDStwoSubj} \\
9 & MiniRocket & \R{StepIDSoneRestBal} & \R{StepIDSoneRestSubj} & \R{StepIDStwoBal} & \R{StepIDStwoSubj} \\
10 & Ensemble & \R{StepJDSoneRestBal} & \R{StepJDSoneRestSubj} & \R{StepJDStwoBal} & \R{StepJDStwoSubj} \\
\hline
\multicolumn{2}{l}{\textbf{Proposed pipeline (Chapter~6)}} & \textbf{\R{MainDSoneRestBal}} & \textbf{\R{MainDSoneRestSubj}} & \textbf{\R{MainDStwoBal}} & \textbf{\R{MainDStwoSubj}} \\
\hline
\end{tabular}\\[2pt]
{\footnotesize \dstwonote}
\end{table}

\begin{figure}[H]\centering
\genfig[0.9\textwidth]{progression}
\caption{DS-2: balanced accuracy on unseen children after each step of the preliminary study and
for the proposed pipeline; the number above each bar is the number of children (of 121)
classified correctly. Dashed line: chance.}
\label{fig:progression}
\end{figure}

\section{What the preliminary study showed}
\begin{enumerate}
\item \textbf{Relations between channels carry the most information about ADHD.} Per-channel
      features (Steps~1, 3 and~5) stay at \R{StepCDStwoBal}\% or below on DS-2, connectivity
      (Step~4) reaches \R{StepDDStwoBal}\%, and the Riemannian tangent space of the covariance
      matrix (Step~7) reaches \R{StepGDStwoBal}\% (\R{StepGDStwoSubj} children), the best of all
      steps. Deep networks (Step~8) and MiniRocket (Step~9) do not exceed it.
\item \textbf{The decomposition alone is not the source of accuracy.} In a ladder that replaced
      SMVMD by fixed bands, VMD and MVMD while keeping the same nine features, the raw signal did
      as well as SMVMD (\R{DStwoLadderRawLr}\% versus \R{DStwoLadderSmvmdLr}\% with logistic
      regression on DS-2). The weakness of the starting method lies in its per-channel
      statistical features, not in the decomposition.
\item \textbf{DS-1 is separable by many methods.} From Step~7 onwards, several methods classify
      all 14 participants correctly, so DS-1 cannot rank these methods.
\end{enumerate}

These findings motivate the proposed pipeline. The covariance of the channels, analysed with
Riemannian geometry, is the strongest representation; the starting method's decomposition can
supply it with frequency-specific components, provided that the window-specific modes are first
made comparable across windows. The pipeline therefore (i) keeps SMVMD, (ii) groups its modes
into common modes, (iii) describes every common mode by a covariance matrix in the Riemannian
tangent space, and (iv) replaces the linear or kernel classifier by gradient-boosted trees.
